\section{Deep Sequence Models on Representative Series}
\label{sec:deep}

This section fits recurrent sequence models (RNN, GRU, LSTM) on the same four representative series used in \cref{sec:eda:repseries} and \cref{sec:classical}. Compared with the midway report, the principal additions are a feature-aware variant of every architecture, bootstrap confidence intervals on every reported skill score, and a direct visual comparison against the classical winners from \cref{tab:classical_leaderboard}. The headline finding is asymmetric: deep models substantively beat classical on the most-volatile series (with statistical significance), badly underperform on the highest-weight series, and roughly tie on the other two.

\subsection{Setup and architectures}
\label{sec:deep:setup}

We reuse the four representative series and the adaptive 80/20 chronological splits from \cref{tab:classical_splits}, so per-series training and validation lengths are identical to those used in the classical study. The metric, weighted skill score, is also identical, scored on the validation segment.

Three architectures are fit with matching shapes:

\begin{itemize}
  \item \textbf{RNN} --- single recurrent layer with hidden size 32 and \texttt{tanh} non-linearity.
  \item \textbf{GRU} --- single gated-recurrent layer with hidden size 32.
  \item \textbf{LSTM} --- single long-short-term-memory layer with hidden size 32.
\end{itemize}

All three are trained for up to 80 epochs with early stopping (patience 10) on a held-out 15\% slice of the per-series \emph{training} segment. The validation segment used for scoring is held out entirely. Optimisation uses Adam at learning rate $10^{-3}$, batch size 32, and MSE loss on standardised one-step windows of length $L = 24$. We do not run a hyperparameter sweep in this lean version; defaults are fixed once and held constant across all (series, architecture, mode) configurations.

\subsection{Two input modes}
\label{sec:deep:modes}

Each architecture is fit twice, once in each of two input modes:

\begin{itemize}
  \item \textbf{target-only} --- the input at each step is the lagged target alone, shape $(B, L, 1)$. This matches the midway setup.
  \item \textbf{feature-aware} --- the input concatenates the lagged target with the lagged top-$k$ features, shape $(B, L, 1 + k)$. We use $k = 10$, picking the features with the largest absolute Pearson correlation against $y_{\text{target}}$ on a 200{,}000-row stratified sample (the same ranking surfaced in \cref{sec:eda:feature_audit}).
\end{itemize}

The recursive multi-step forecast at validation time uses the predicted $y$ at each step as the next-step's lagged target. Future feature values, when used, are taken from the dataset; they are exogenous knowns at validation time, not labels. The open-loop one-step diagnostic is computed by always feeding the true previous $y$ at each step --- analogous to the protocol introduced in \cref{sec:classical:protocol}.

\subsection{Per-series leaderboard}
\label{sec:deep:leaderboard}

\Cref{tab:deep_leaderboard} reports the best deep configuration per series under recursive-multi-step weighted skill, with 95\% bootstrap intervals from 500 resamples of the validation rows. The corresponding best classical model (from \cref{tab:classical_leaderboard}) and the Naive (lag-1) floor are shown alongside.

\begin{table}[h]
\centering
\caption{Per-series deep leaderboard. ``Best deep'' shows the winning architecture and input mode; ``Skill (95\% CI)'' is the recursive-multi-step weighted skill of that configuration with a 500-sample bootstrap interval. ``Classical'' is the corresponding winner from \cref{tab:classical_leaderboard}.}
\label{tab:deep_leaderboard}
\begin{tabular}{l l c c c}
\toprule
Series & Best deep & Skill (95\% CI) & Naive floor & Classical \\
\midrule
Longest history     & GRU / feature-aware  & $0.04$ \,$[0.00, 0.61]$ & $0.00$ & AR$(3,0,0)$ $0.04$ \\
Highest total weight & RNN / target-only    & $0.18$ \,$[0.00, 0.33]$ & $0.44$ & SES $0.47$ \\
Most volatile       & GRU / feature-aware  & $\mathbf{0.87}$ \,$[\mathbf{0.83}, \mathbf{0.90}]$ & $0.67$ & Rolling Mean $0.80$ \\
Most stable         & GRU / target-only    & $0.15$ \,$[0.00, 0.23]$ & $0.00$ & Rolling Mean $0.17$ \\
\bottomrule
\end{tabular}
\end{table}

The most-volatile series is the only case with a statistically significant deep-over-classical improvement: the bootstrap CI $[0.83, 0.90]$ excludes the classical Rolling Mean point estimate $0.80$. The feature-aware GRU produced this win; the target-only variant of the same architecture scored only $0.32$, so the features themselves are doing meaningful work on this series.

The highest-weight result is the cleanest negative finding. The best deep configuration scores $0.18$, well below the SES classical winner at $0.47$ and even below the Naive floor at $0.44$. Recursive multi-step recurrence on a series operating at scale $10^{-4}$ accumulates error faster than smoothing absorbs it; with only 124 training rows, the network has too little data to lock in the mean dynamics that SES captures trivially.

The longest-history and most-stable cases are inside-noise ties with classical. The wide CIs on both --- the longest-history CI runs all the way from $0.00$ to $0.61$ --- reflect a 32--43-row validation segment where every model's score is bootstrap-fragile.

\Cref{fig:deep_skill_bars} shows the same data as a per-series bar chart, with the Naive and best-classical reference lines drawn. The within-series ranking is what matters; cross-series mean skills would be dominated by the volatile case.

\begin{figure}[h]
\centering
\includegraphics[width=\linewidth]{02_per_series_skill_bars}
\caption{Recursive-multi-step weighted skill for every (architecture, input mode) configuration on each series. Bar colour encodes architecture (RNN green, GRU red, LSTM purple); bar opacity encodes mode (solid for feature-aware, faded for target-only). The Naive (lag-1) floor and the best classical winner from \cref{tab:classical_leaderboard} are shown as dashed vertical lines.}
\label{fig:deep_skill_bars}
\end{figure}

\subsection{Deep versus classical forecast overlays}
\label{sec:deep:overlay}

\Cref{fig:deep_overlay} shows the per-series train and validation segments alongside the best deep configuration's recursive-multi-step forecast, the corresponding best classical model's forecast (from \cref{sec:classical:overlay}), and the Naive (lag-1) floor. Skill annotations on each panel make the deep-versus-classical comparison directly readable.

\begin{figure}[h]
\centering
\includegraphics[width=\linewidth]{02_forecast_overlay_2x2}
\caption{Forecast overlays for the four representative series. Black: training segment. Orange: validation truth. Brown: best classical winner from \cref{tab:classical_leaderboard}. Coloured (architecture-specific): best deep configuration. Grey dashed: Naive (lag-1) floor. Vertical dashed line: per-series train/validation cutoff. Skill annotations are weighted skill scores against the validation segment.}
\label{fig:deep_overlay}
\end{figure}

The most-volatile panel is the only case where the deep forecast tracks the validation trajectory meaningfully better than the classical Rolling Mean (which is flat by construction). On the highest-weight panel, the deep forecast curves away from the truth more aggressively than the classical SES line, consistent with the recursive-error-compounding explanation above. The longest-history and most-stable panels show every method clustered near zero --- the empirical content of the ``no signal to extract'' framing carried over from \cref{sec:classical:conclusions}.

\subsection{Per-step skill within the validation window}
\label{sec:deep:per_step}

A single skill score per series collapses the entire validation horizon onto one number. Computing weighted skill at horizons $h \in \{1, 5, 10, 25\}$ within the validation window separates short-horizon performance (where one-step memory is what matters) from long-horizon performance (where compounding errors dominate). The training-curve and per-step-skill diagnostics are produced by the notebook (\texttt{02\_per\_step\_skill.pdf} and \texttt{02\_training\_curves.pdf}) and are not embedded here for brevity. Two qualitative observations:

\begin{itemize}
  \item On the most-volatile series, GRU / feature-aware retains skill close to $0.85$ at every horizon, indicating the model has captured a stable feature--target relationship rather than an early-window fluke.
  \item On the highest-weight series, the open-loop one-step skill ($0.66$ for LSTM) is substantially higher than the recursive-multi-step skill ($0.00$). The gap is the entirety of the deep failure on that series: feedback at inference is what classical SES is implicitly providing through its update rule, and recurrent models without that feedback cannot match it.
\end{itemize}

This observation motivates the next chapter. If recurrent models with feedback at inference could be made deployable, they might recover the highest-weight loss; the H1-aggregation experiment in \cref{sec:cumulation} explores a structurally different way to inject feedback by exploiting the multi-horizon coupling property of the panel.

\subsection{What this study can and cannot say}
\label{sec:deep:conclusions}

Three points to carry forward:

\begin{enumerate}
  \item \textbf{Deep sequence models help where classical fails for sequence-structural reasons, and not elsewhere.} The most-volatile series has dynamic structure that smoothing cannot capture; the feature-aware GRU captures it. The highest-weight series has structure that smoothing captures trivially via its update rule; recursive recurrence cannot match that without inference-time feedback. The other two series have no exploitable signal at this scale.
  \item \textbf{The feature-aware variant is not a free win.} It produced the volatile-series victory but did not improve the highest-weight result. With only ${\sim}120$ training rows and 10 added feature channels, capacity outpaces data on series where lagged target alone already saturates the available signal.
  \item \textbf{Per-series fitting still bounds the result, exactly as in \cref{sec:classical:conclusions}.} Each model sees ${\sim}120$ training rows and zero cross-series structure. The next chapter (global LightGBM) lifts that limit; the H1-aggregation experiment after that uses the multi-horizon coupling in a way no per-series model can.
\end{enumerate}
